\begin{textsample}{POS Dim 3 – human – Score 6.00 – mg/mg\_ef\_linguagens\_ingles\_1\_p1.txt}  \label{ex:f3_pos_011}
O Currículo Referência de Minas Gerais do componente \textbf{curricular} Língua Inglesa é o resultado de um diálogo entre os Conteúdos Básicos \textbf{Comuns} ( CBC ) já utilizado nas escolas estaduais e a Base \textbf{Nacional} Comum \textbf{Curricular} ( BNCC ) homologada em dezembro de 2017 pelo Conselho Nacional de Educação ( CNE ).
No \textbf{documento}, são apresentadas ao professor as expectativas de aprendizagem representadas pelas habilidades listadas por ano de escolaridade na Língua Inglesa.
As orientações metodológicas e os referenciais teóricos são de grande importância, quando devidamente apropriadas pelo educador são ferramentas poderosas para qualificar o trabalho docente, adicionando infinitas possibilidades de reflexão sobre o fazer pedagógico.
O conhecimento da língua como sistema de regras gramaticais sempre esteve em \textbf{primeiro} plano e os conteúdos relacionados à descrição da estrutura sintática da língua constituíam os eixos organizativos do Currículo, confinando o estudo do léxico a mero objeto para o preenchimento de lacunas das estruturas estudadas.

% matched lemmas: comum, curricular, documento, nacional, primeiro
\end{textsample}
